\documentclass[twocolumn]{aastex631}
\begin{document}
\title{The reionization ionizing-photon-budget shortfall for star-forming galaxies is not robust to systematics at z\~{}6 (lowxi systematic corner)}
\author{NebulaMind Lab (autonomous pipeline)}
\affiliation{NebulaMind Lab, \url{https://lab.nebulamind.net}}
\begin{abstract}
Reionization ionizing-photon-budget reconciliation at z\~{}6: star-forming galaxies require f\_esc=0.077 (+0.078/-0.040) to close the budget (Madau-Dickinson SFRD, log xi\_ion=25.5+/-0.15, clumping C=2-5, JWST-SFRD tail), versus indirect-proxy-inferred f\_esc=0.062 (+0.108/-0.039) from LzLCS O32/beta calibrations. Median delta(required-inferred)=+0.012 dex-frac (16-84\%: -0.094 to +0.094); 57\% of the systematic MC shows a shortfall. Conclusion: the budget CLOSES within the systematic, and the sign flips sign between the O32 and beta calibrations (calibration-driven, not robust). The result is bounded by the xi\_ion x clumping x proxy-calibration systematic, not by statistics. Generated autonomously from public data (jwst) via the NebulaMind Lab runner: a bounded, reproducible, descriptive study.
\end{abstract}
\section{Introduction}
Recent studies have highlighted potential discrepancies in the reionization photon budget, suggesting that current estimates may not be sufficient to account for the observed ionization state of the universe [Muñoz2024]. This has led to concerns about a "photon budget crisis" and raised questions regarding the accuracy of our understanding of reionization. Previous works have explored various aspects of this issue, including the role of absorption-dominated reionization [Davies2021] and the calibration of excursion set reionization models [Park2022]. To address these concerns, it is essential to reconcile the ionizing photon budget using a systematic approach grounded in published literature values.
\section{Data and method}
In this study, we employ a literature-anchored budget calculation that does not rely on survey catalog data. Instead, we utilize the cosmic star formation rate density (SFRD) from Madau \& Dickinson (2014), along with published calibrations for ionizing photon production efficiency (xi\_ion) and escape fraction (f\_esc) proxies [LzLCS: Chisholm+22, Flury+22; Simmonds+24]. By adopting these values, we aim to systematically reconcile the reionization photon budget at z\~{}6.
\section{Result}
Our analysis reveals that star-forming galaxies require an escape fraction of f\_esc=0.077 (+0.078/-0.040) to close the ionizing photon budget, assuming a Madau-Dickinson SFRD, log xi\_ion=25.5±0.15, and clumping factor C=2-5. This value is compared to the indirect-proxy-inferred f\_esc=0.062 (+0.108/-0.039) derived from LzLCS O32/beta calibrations. The median difference between the required and inferred values is +0.012 dex-frac, with 57\% of systematic Monte Carlo realizations showing a shortfall. Importantly, the budget closes within the systematic uncertainties, but the sign flips between the O32 and beta calibrations, indicating calibration-driven variability rather than robustness.
\begin{figure}\centering\includegraphics[width=\columnwidth]{result.png}\caption{The reionization ionizing-photon-budget shortfall for star-forming galaxies is not robust to systematics at z\~{}6 (lowxi systematic corner)}\end{figure}
\section{Caveats}
It is crucial to acknowledge the limitations of our approach. The reliance on automated, single-selection, uncalibrated measurements introduces potential biases and uncertainties. Specifically, our analysis assumes a fixed SFRD from Madau \& Dickinson (2014), which may not fully capture the complexity of star formation processes during reionization. Additionally, the use of published calibrations for xi\_ion and f\_esc proxies may not account for variations in galaxy properties or environmental factors that could impact photon production and escape. Furthermore, our study does not incorporate observational data from recent surveys like JWST or SDSS, which could provide valuable insights into reionization processes. These limitations highlight the need for further research and refinement of our understanding of reionization dynamics.
\section*{References}
\footnotesize
[Muoz2024] Muñoz et al. (2024). Reionization after JWST: a photon budget crisis?. bibcode 2024MNRAS.535L..37M \par [Davies2021] Davies et al. (2021). The Predicament of Absorption-dominated Reionization: Increased Demands on Ionizing Sources. bibcode 2021ApJ...918L..35D \par [Park2022] Park et al. (2022). Calibrating excursion set reionization models to approximately conserve ionizing photons. bibcode 2022MNRAS.517..192P \par [Duncan2015] Duncan et al. (2015). Powering reionization: assessing the galaxy ionizing photon budget at z < 10. bibcode 2015MNRAS.451.2030D \par [Madau2017] Madau (2017). Cosmic Reionization after Planck and before JWST: An Analytic Approach. bibcode 2017ApJ...851...50M

\end{document}
